\section{Алгоритмы сортировки}

\subsection{Постановка задачи и свойства сортировок}

\textbf{Сортировка} --- перестановка элементов последовательности в соответствии с заданным отношением порядка. Для ключей $k_i$ требуется получить
\[
k'_1\le k'_2\le\cdots\le k'_n.
\]
Если сортируются записи $(k_i,d_i)$, сравнение выполняется по ключу, а связанные данные переносятся вместе с ним.

Основные характеристики алгоритма:
\begin{itemize}
\item время и дополнительная память;
\item \textbf{устойчивость}: элементы с одинаковыми ключами сохраняют исходный взаимный порядок;
\item работа \textbf{на месте} (in-place) --- практически без дополнительного массива;
\item адаптивность к уже частично отсортированным данным;
\item использование только сравнений или дополнительных свойств ключа.
\end{itemize}

\subsection{Простые квадратичные алгоритмы}

\textbf{Insertion Sort}. Уже отсортирован префикс $a_1,\ldots,a_{i-1}$; очередной элемент вставляется в нужное место со сдвигом больших элементов.
\[
T_{\rm worst}(n)=\Theta(n^2),
\qquad
T_{\rm best}(n)=\Theta(n).
\]
Алгоритм устойчив, работает на месте и эффективен для небольших или почти упорядоченных массивов.

\textbf{Selection Sort}. На каждом шаге выбирается минимум оставшейся части и переносится в начало:
\[
T(n)=\Theta(n^2)
\]
для любого расположения данных; обменов всего $O(n)$. Обычная реализация неустойчива.

\textbf{Bubble Sort}. Соседние элементы многократно меняются местами при нарушении порядка. Худший и средний случаи $\Theta(n^2)$; с признаком отсутствия обменов лучший случай $\Theta(n)$. Алгоритм устойчив, но практически редко эффективен.

\subsection[Сортировки порядка n log n сравнениями]{Сортировки $\Theta(n\log n)$ сравнениями}

\textbf{Merge Sort} делит массив пополам, сортирует части и за линейное время сливает две отсортированные последовательности:
\[
T(n)=2T(n/2)+\Theta(n)=\Theta(n\log n).
\]
Эта оценка верна и в худшем случае. Классическая реализация устойчива и требует $O(n)$ дополнительной памяти; особенно удобна для внешней сортировки.

\textbf{Quick Sort} выбирает опорный элемент и разделяет данные относительно него. При сбалансированном разбиении
\[
T(n)=\Theta(n\log n),
\]
а в худшем случае
\[
T(n)=\Theta(n^2).
\]
При случайном выборе опорного элемента ожидаемое время равно
\[\Theta(n\log n).\]
Классическая реализация на месте неустойчива; ожидаемая глубина рекурсии равна $O(\log n)$, а в худшем случае --- $O(n)$.

\textbf{Heap Sort} использует двоичную максимальную кучу. Построение кучи занимает $O(n)$, затем $n$ извлечений максимума требуют по $O(\log n)$:
\[
T(n)=\Theta(n\log n)
\]
в худшем случае. Алгоритм работает на месте и неустойчив.

\subsection{Сортировки без сравнений}

Нижняя граница $\Omega(n\log n)$ относится только к сортировкам, получающим информацию о порядке исключительно из сравнений.

\textbf{Counting Sort}. Для целых ключей из диапазона $0,\ldots,k-1$ подсчитывается число элементов каждого ключа. Стандартная устойчивая реализация имеет
\[
T(n,k)=O(n+k),
\qquad
M(n,k)=O(n+k)
\]
(если нужен только подсчёт частот без устойчивого массива результата, дополнительная память может быть $O(k)$).

\textbf{Radix Sort} последовательно сортирует ключи по разрядам, используя устойчивую внутреннюю сортировку. При $d$ разрядах:
\[
T=O\bigl(d(n+k)\bigr).
\]

\textbf{Bucket Sort} распределяет элементы по ``карманам'' и сортирует их независимо. При благоприятном, например близком к равномерному, распределении возможно ожидаемое $O(n)$, но худший случай остаётся квадратичным.

\subsection{Нижняя оценка для сортировок сравнениями}

\textbf{Теорема.} Любой алгоритм сортировки $n$ различных элементов, использующий только попарные сравнения, в худшем случае выполняет
\[
\Omega(n\log n)
\]
сравнений.

\textbf{Доказательство через дерево решений.} Каждое сравнение даёт ветвление, а для $n$ различных элементов необходимо различить $n!$ возможных перестановок. Поэтому дерево решений имеет как минимум $n!$ листьев. Если его высота равна $h$, то
\[
2^h\ge n!,
\qquad
h\ge\log_2(n!).
\]
По формуле Стирлинга
\[
\log(n!)=\Theta(n\log n),
\]
что и даёт нижнюю оценку. Следовательно, Merge Sort и Heap Sort асимптотически оптимальны среди сортировок сравнениями по худшему времени.

\subsection{Сравнение алгоритмов}

\begin{center}
\begin{tabular}{|l|c|c|c|c|}
\hline
\textbf{Алгоритм} & \textbf{лучший} & \textbf{средний} & \textbf{худший} & \textbf{память} \\
\hline
Insertion & $O(n)$ & $O(n^2)$ & $O(n^2)$ & $O(1)$ \\
Selection & $O(n^2)$ & $O(n^2)$ & $O(n^2)$ & $O(1)$ \\
Merge & $O(n\log n)$ & $O(n\log n)$ & $O(n\log n)$ & $O(n)$ \\
Quick & $O(n\log n)$ & $O(n\log n)$ & $O(n^2)$ & $O(\log n)$ ожидаемо \\
Heap & $O(n\log n)$ & $O(n\log n)$ & $O(n\log n)$ & $O(1)$ \\
Counting & $O(n+k)$ & $O(n+k)$ & $O(n+k)$ & $O(n+k)$* \\
\hline
\end{tabular}
\end{center}

\noindent *Для стандартной устойчивой реализации с выходным массивом.

Если данные не помещаются в оперативную память, существенной становится стоимость ввода-вывода. Классический вариант --- \textbf{внешняя сортировка слиянием}: в памяти сортируются отдельные серии (runs), после чего они многопутево сливаются последовательным чтением больших блоков.


\subsection{Корректность простых сортировок через инвариант}

Для сортировки вставками удобен инвариант цикла: перед началом шага $i$
подмассив $a_1,\ldots,a_{i-1}$ уже отсортирован и содержит те же элементы, что и
исходный префикс. Вставка $a_i$ сохраняет это свойство для префикса длины $i$.
После последнего шага инвариант распространяется на весь массив.

Подобные инварианты позволяют доказывать корректность независимо от оценки
сложности. Для selection sort после $i$ шагов первые $i$ позиций содержат $i$
наименьших элементов в правильном порядке относительно выбранной схемы.

\subsection{Partition в быстрой сортировке}

Ключевой этап Quick Sort --- разбиение массива относительно опорного элемента
$p$. После partition должны выполняться свойства
\[
a_i\le p\quad\text{слева},\qquad
a_j\ge p\quad\text{справа}
\]
(точные неравенства зависят от схемы Hoare/Lomuto). После этого части можно
сортировать независимо.

Худший случай возникает при крайне несбалансированных разбиениях:
\[
T(n)=T(n-1)+\Theta(n)=\Theta(n^2).
\]
При случайном выборе pivot ожидаемое число сравнений имеет порядок
$\Theta(n\log n)$, а ожидаемая глубина рекурсии --- $O(\log n)$; при этом
худшая глубина по-прежнему может достигать $O(n)$. Поэтому рандомизация защищает
от регулярно неудачных структур входа в среднем, но не устраняет худший случай.

\subsection{Куча и Heap Sort}

Максимальная двоичная куча удовлетворяет
\[
a_{parent(i)}\ge a_i.
\]
Построение кучи снизу алгоритмом heapify занимает $O(n)$, хотя последовательная
вставка $n$ элементов дала бы $O(n\log n)$. Затем максимум находится в корне,
переставляется в конец массива, размер кучи уменьшается, а свойство кучи
восстанавливается за $O(\log n)$.

Отсюда гарантированное
\[
T(n)=\Theta(n\log n),\qquad M(n)=O(1)
\]
для массивной реализации.

\subsection{Практический выбор алгоритма}

Асимптотика --- не единственный критерий:
\begin{itemize}
    \item insertion sort эффективен на коротких или почти отсортированных блоках;
    \item merge sort удобен при требовании устойчивости и для внешней сортировки;
    \item quicksort имеет хорошую локальность памяти и малые константы;
    \item heapsort даёт гарантированный $O(n\log n)$ без дополнительного массива;
    \item counting/radix используют структуру ключа и могут быть быстрее
    сравнительной нижней границы.
\end{itemize}

Поэтому библиотечные алгоритмы нередко гибридны: рекурсивная сортировка
переключается на insertion sort на малых подмассивах, а introsort начинает как
quicksort и при слишком большой глубине переключается на heapsort, сохраняя
хорошую среднюю скорость и гарантированный худший $O(n\log n)$.

\subsection{Внешняя сортировка}

Если данные не помещаются в оперативную память, стоимость определяется прежде
всего вводом-выводом. Внешняя сортировка слиянием сначала формирует в памяти
отсортированные серии, записывает их на диск, затем выполняет многопутевое
слияние. Большие последовательные операции чтения/записи существенно выгоднее
случайного доступа к диску.


\subsection{Что важно сказать на экзамене}

Нужно не перечислять десятки алгоритмов, а уверенно сравнить три основных $O(n\log n)$-метода:
\begin{itemize}
\item Merge Sort --- устойчивость и гарантированное $O(n\log n)$;
\item Quick Sort --- высокая практическая скорость и $O(n\log n)$ в среднем;
\item Heap Sort --- гарантированное $O(n\log n)$ при $O(1)$ дополнительной памяти.
\end{itemize}
Затем сформулировать нижнюю оценку $\Omega(n\log n)$ для сортировок сравнениями и объяснить, почему Counting/Radix её не нарушают.

\newpage
